\chapter{Weight separation, integrality, and independence of \texorpdfstring{\(\ell\)}{ell}}
\label{chap:integrality}

The zeta function has integral Taylor coefficients for the elementary
reason that it is an Euler product over closed points.  Its
cohomological factors were constructed over \(\Ql\), and their roots
initially live only in \(\Qbarl\).  This chapter proves, in the required
noncircular order, that those roots are algebraic over \(\mathbf Q\), that
the individual factors lie in \(\mathbf Z[T]\), and that they are
independent of \(\ell\):
\[
\begin{split}
 &\text{purity for every }\iota:\Qbarl\simeq\mathbf C\\
 &\quad\Longrightarrow\text{no cancellation between distinct weights}\\
 &\quad\Longrightarrow\text{roots occur in the reduced integral
 numerator or denominator of }Z\\
 &\quad\Longrightarrow\text{algebraicity over }\mathbf Q
 \text{ and Galois-stable weight pieces}\\
 &\quad\Longrightarrow\mathbf Q\text{-rational factors}
 \Longrightarrow\mathbf Z\text{-integral factors and }\ell\text{-independence}.
\end{split}
\tag{14.1}\label{eq:arithmetic-order}
\]
In particular, algebraicity over \(\mathbf Q\) is a conclusion in the
middle of the argument, not part of the initial definition of weight.

\section{Universal realization-purity}

\begin{definition}
Let \(\alpha\in\Qbarl^\times\), let \(q>1\), and let \(w\in\mathbf Z\).
We say that \(\alpha\) is \emph{universally \(\iota\)-pure of weight
\(w\)} if
\[
  |\iota(\alpha)|=q^{w/2}
\]
for every abstract field isomorphism
\(\iota\colon\Qbarl\xrightarrow{\sim}\mathbf C\).
\end{definition}

The two algebraically closed fields in this definition have
characteristic \(0\), cardinality \(2^{\aleph_0}\), and transcendence
degree \(2^{\aleph_0}\) over \(\mathbf Q\), so such isomorphisms exist.
No topology is imposed on them; a continuous map from the \(\ell\)-adic
topology to the complex topology would be the wrong notion.

\begin{lemma}[Invariance and disjointness]
\label{lem:universal-weight-invariance}
Let \(\alpha\in\Qbarl^\times\).
\begin{enumerate}
\item If \(\alpha\) is universally \(\iota\)-pure of weight \(w\) and
      \(\sigma\in\operatorname{Aut}(\Qbarl/\mathbf Q)\), then
      \(\sigma(\alpha)\) is universally \(\iota\)-pure of weight \(w\).
\item The same nonzero element cannot be universally
      \(\iota\)-pure of two distinct weights.
\end{enumerate}
\end{lemma}

\begin{proof}
For every \(\iota\colon\Qbarl\simeq\mathbf C\), the composite
\(\iota\circ\sigma\) is another such isomorphism.  Hence
\[
 |\iota(\sigma\alpha)|
 =|(\iota\circ\sigma)(\alpha)|=q^{w/2}.
\]
This proves the first assertion.  If \(\alpha\) had weights \(w\) and
\(w'\), evaluation under any one isomorphism would give
\(q^{w/2}=q^{w'/2}\).  Since \(q>1\), this implies \(w=w'\).
\end{proof}

\begin{lemma}[Extension and fixed fields]
\label{lem:qbarl-automorphism-fixed-field}
The following field-theoretic facts hold.
\begin{enumerate}
\item Every \(\mathbf Q\)-embedding of a finite algebraic subextension
      \(E\subset\Qbarl\) into \(\mathbf C\) extends to an isomorphism
      \(\Qbarl\xrightarrow{\sim}\mathbf C\).
\item Every \(\mathbf Q\)-embedding between finite algebraic
      subextensions of \(\Qbarl\) extends to an element of
      \(\operatorname{Aut}(\Qbarl/\mathbf Q)\).
\item The fixed field of
      \(\operatorname{Aut}(\Qbarl/\mathbf Q)\) is \(\mathbf Q\).
\end{enumerate}
\end{lemma}

\begin{proof}
For (1), first extend the given embedding over the relative algebraic
closure of \(\mathbf Q\) in \(\Qbarl\).  Choose transcendence bases on
the source and target complements.  Both bases have cardinality
\(2^{\aleph_0}\), so a bijection between them, together with the
algebraic embedding already chosen, gives an isomorphism of the purely
transcendental subfields.  Since source and target are algebraically
closed, the isomorphism extends across their algebraic closures.  The
same argument with target \(\Qbarl\) proves (2).

For (3), let \(x\in\Qbarl\setminus\mathbf Q\).  If \(x\) is algebraic
over \(\mathbf Q\), then either it has a distinct conjugate, in which case
(2) moves it, or its minimal polynomial is linear, which would put it in
\(\mathbf Q\).  If \(x\) is transcendental, include it in a
transcendence basis and permute that basis; the permutation extends to an
automorphism of the algebraic closure and moves \(x\).  Thus no element
outside \(\mathbf Q\) is fixed by every automorphism.
\end{proof}

\section{No cancellation before algebraicity}

Fix an auxiliary prime \(\ell\), and set
\[
  O_\ell(T)=\prod_{\substack{0\leq i\leq2d\\ i\ {\rm odd}}}
       P_{i,\ell}(T),\qquad
  E_\ell(T)=\prod_{\substack{0\leq i\leq2d\\ i\ {\rm even}}}
       P_{i,\ell}(T).
\]
These polynomials lie in \(1+T\Ql[T]\), and
\[
  Z(X_0,T)=\frac{O_\ell(T)}{E_\ell(T)}.
  \tag{14.2}\label{eq:odd-even-zeta-final}
\]

\begin{lemma}[Cohomological numerator and denominator are coprime]
\label{lem:odd-even-coprime-before-Q}
The polynomials \(O_\ell\) and \(E_\ell\) are coprime in
\(\Qbarl[T]\).  More generally, \(P_{i,\ell}\) and \(P_{j,\ell}\) have
no common root when \(i\neq j\).
\end{lemma}

\begin{proof}
Suppose \(P_{i,\ell}\) and \(P_{j,\ell}\) shared the ordinary root
\(\rho\).  Frobenius is invertible, so \(\rho\neq0\), and
\(\alpha=\rho^{-1}\) would be a reciprocal root of both polynomials.
By \cref{prop:all-realization-purity}, it would be universally
\(\iota\)-pure of weights \(i\) and \(j\).  The disjointness part of
\cref{lem:universal-weight-invariance} forces \(i=j\).  In particular an
odd factor and an even factor cannot meet, proving the first assertion.
\end{proof}

Now use integrality of the \emph{whole} zeta series.  By
\cref{prop:euler-product,thm:cohomological-rationality},
\[
 Z(X_0,T)\in1+T\mathbf Z[[T]]\cap\mathbf Q(T).
\]
The reduced integral Fatou presentation
\cref{cor:reduced-integral-presentation} gives coprime
\(A,B\in\mathbf Z[T]\) with
\[
  Z(X_0,T)=A(T)/B(T),\qquad A(0)=B(0)=1.
  \tag{14.3}\label{eq:reduced-integral-zeta}
\]

\begin{proposition}[Parity products are the integral reduced factors]
\label{prop:parity-products-integral}
For every \(\ell\),
\[
  O_\ell(T)=A(T),\qquad E_\ell(T)=B(T)
\]
inside \(\Qbarl[T]\).
\end{proposition}

\begin{proof}
The fraction \(O_\ell/E_\ell\) is reduced over the algebraically closed
field \(\Qbarl\) by
\cref{lem:odd-even-coprime-before-Q}.  The fraction \(A/B\) remains
reduced after scalar extension from \(\mathbf Q\), since Bezout
coprimality is preserved.  Reduced presentations over a field are unique
up to a common nonzero scalar.  All four polynomials have constant term
\(1\), so that scalar is \(1\).
\end{proof}

\begin{corollary}[Arithmetic algebraicity of eigenvalues]
\label{cor:eigenvalues-algebraic-over-Q}
Every Frobenius eigenvalue on every \(V_{i,\ell}\) is algebraic over
\(\mathbf Q\).
\end{corollary}

\begin{proof}
Every reciprocal root of every \(P_{i,\ell}\) occurs in the product
\(O_\ell\) or \(E_\ell\).  By
\cref{prop:parity-products-integral}, it is therefore a reciprocal root
of the polynomial \(A\) or \(B\) with rational coefficients.  Its
inverse is algebraic over \(\mathbf Q\).
\end{proof}

This is the precise point where it becomes legitimate to regard the
eigenvalues as algebraic numbers.  Nothing in the proof of no
cancellation used that conclusion.

\section{Recovery of the rational weight factors}

The roots of \(A\) and \(B\), viewed inside \(\Qbarl\), are algebraic over
\(\mathbf Q\).  Partition their reciprocal roots by universal
\(\iota\)-weight.  By
\cref{lem:universal-weight-invariance}, every part is stable under
\(\operatorname{Aut}(\Qbarl/\mathbf Q)\), and distinct parts are
disjoint.  The part of weight \(i\), with multiplicity, is exactly the
reciprocal-root multiset of \(P_{i,\ell}\).

\begin{theorem}[Rationality of individual factors]
\label{thm:rational-individual-factors}
For every \(i\) and every \(\ell\neq\operatorname{char}(\Fq)\),
\[
  P_{i,\ell}(T)\in\mathbf Q[T].
\]
Moreover, if \(\alpha\) is one of its reciprocal roots, then for every
embedding
\(\tau\colon\mathbf Q(\alpha)\hookrightarrow\mathbf C\),
\[
  |\tau(\alpha)|=q^{i/2}.
  \tag{14.4}\label{eq:classical-all-embeddings}
\]
\end{theorem}

\begin{proof}
An automorphism of \(\Qbarl\) over \(\mathbf Q\) permutes the roots of
\(A\) and \(B\), because those polynomials have rational coefficients.
It preserves each universal-weight part by
\cref{lem:universal-weight-invariance}.  Hence it preserves the
reciprocal-root multiset of \(P_{i,\ell}\).  Since
\(P_{i,\ell}(0)=1\), the polynomial itself is fixed coefficientwise by
every such automorphism.  The fixed-field assertion of
\cref{lem:qbarl-automorphism-fixed-field} gives
\(P_{i,\ell}\in\mathbf Q[T]\).

Let \(\tau\) be any complex embedding of the now-known number field
\(\mathbf Q(\alpha)\).  By
\cref{lem:qbarl-automorphism-fixed-field}(1), it extends to an isomorphism
\(\iota\colon\Qbarl\simeq\mathbf C\).  Universal realization-purity then
gives
\[
 |\tau(\alpha)|=|\iota(\alpha)|=q^{i/2}.
\]
Thus the universal \(\ell\)-adic formulation has yielded the classical
all-number-field-embeddings formulation.
\end{proof}

\section{Gauss's lemma and integral factors}

\begin{theorem}[Integrality of the individual factors]
\label{thm:integral-frobenius-factors}
For every \(i\) and every auxiliary prime \(\ell\),
\[
  P_{i,\ell}(T)\in\mathbf Z[T].
\]
Consequently every Frobenius eigenvalue is an algebraic integer.
\end{theorem}

\begin{proof}
According to the parity of \(i\), the rational polynomial
\(P_{i,\ell}\) is a factor of \(A=O_\ell\) or \(B=E_\ell\).  Its
complement is the product of the remaining factors of the same parity.
Both factor and complement have rational coefficients by
\cref{thm:rational-individual-factors} and constant term \(1\).  The
normalized Gauss lemma \cref{lem:normalized-gauss-factor} therefore
places both in \(\mathbf Z[T]\).

Write
\[
 P_{i,\ell}(T)=1+c_1T+\cdots+c_bT^b
              =\prod_{j=1}^b(1-\alpha_jT).
\]
Each \(\alpha_j\) is a root of the monic integral polynomial
\[
 U^b+c_1U^{b-1}+\cdots+c_b.
\]
It is therefore an algebraic integer.
\end{proof}

\begin{remark}
The Fatou lemma alone provides integral \(A\) and \(B\), but does not
identify their cohomological subfactors.  Universal purity first proves
that the cohomological parity fraction is reduced; only then does
\cref{prop:parity-products-integral} identify it with \(A/B\).
Automorphism-stable weight pieces then give rationality, after which
Gauss's lemma applies.  This order contains no hidden use of individual
integrality.
\end{remark}

\section{Independence of the auxiliary prime}

\begin{theorem}[\(\ell\)-independence]
\label{thm:ell-independence}
For each \(i\), the polynomial \(P_{i,\ell}(T)\) is independent of
\(\ell\neq\operatorname{char}(\Fq)\).  There is therefore a unique
\[
  P_i(T)\in\mathbf Z[T]
\]
such that
\[
  P_i(T)=\det(1-TF\mid H^i_{\et}(X,\Ql))
\]
for every auxiliary prime.  In particular
\[
  b_i=\deg P_i=\dim_{\Ql}H^i_{\et}(X,\Ql)
\]
is independent of \(\ell\).
\end{theorem}

\begin{proof}
The reduced integral polynomials \(A\) and \(B\) in
\eqref{eq:reduced-integral-zeta} are determined by the point counts and
are independent of \(\ell\).  By
\cref{prop:parity-products-integral}, their reciprocal-root multisets are
the union of the odd, respectively even, Frobenius eigenvalue multisets
for every auxiliary prime.

Now work with the roots of the fixed polynomials \(A\) and \(B\) in
\(\Qbar\).  By \cref{thm:rational-individual-factors}, each root has a
classical, all-complex-embeddings weight, and
\cref{lem:universal-weight-invariance} shows distinct weights cannot
meet.  Thus the normalized factor whose reciprocal roots have weight
\(i\) is uniquely determined by \(A/B\), \(q\), and \(i\), with no
reference to \(\ell\).  For every \(\ell\), that unique factor is
\(P_{i,\ell}\).  Hence the polynomials and their degrees agree.
\end{proof}

\begin{corollary}[Betti comparison when a lift exists]
\label{cor:betti-comparison-lift}
Suppose \(X_0\) is the special fiber of a smooth proper scheme over a
mixed-characteristic discrete valuation ring, and let
\(X_{\mathbf C}\) be a complex fiber obtained from an embedding of the
fraction field into \(\mathbf C\).  Then
\[
 b_i=\dim_{\mathbf Q}H^i_{\mathrm{sing}}
       (X_{\mathbf C}(\mathbf C),\mathbf Q).
\]
\end{corollary}

\begin{proof}
Smooth proper base change identifies the etale cohomology of geometric
special and generic fibers.  The characteristic-zero comparison theorem
identifies the latter, after tensoring with \(\Ql\), with the singular
cohomology of the complex analytic fiber.  Taking dimensions proves the
claim.
\end{proof}

\section{Completion of the proof}

\begin{theorem}[The classical Weil conjectures]
\label{thm:weil-conjectures-proved}
All assertions of \cref{thm:weil-conjectures-statement} hold.
\end{theorem}

\begin{proof}
Cohomological rationality is
\cref{thm:cohomological-rationality}; the endpoint factors are
\cref{prop:endpoint-factors}; and the functional equation is
\cref{thm:zeta-functional-equation}.  Universal smooth-proper purity,
followed by \cref{cor:eigenvalues-algebraic-over-Q} and
\cref{thm:rational-individual-factors}, gives the all-complex-embeddings
Riemann-hypothesis statement for algebraic numbers.
\Cref{thm:integral-frobenius-factors} adds algebraic integrality, and
\cref{thm:ell-independence} gives common polynomials and Betti numbers.
These are precisely the conclusions in
\cref{thm:weil-conjectures-statement}.
\end{proof}

\begin{formalization}
The formal development should keep three predicates distinct:
\begin{enumerate}
\item `UniversallyRealizationPure alpha w`, quantified over all field
      isomorphisms \(\Qbarl\simeq\mathbf C\);
\item `IsAlgebraic Q alpha`, proved only after
      \cref{prop:parity-products-integral};
\item `IsIntegral Q alpha`, proved only after the normalized Gauss
      factor theorem.
\end{enumerate}
The reusable arithmetic pipeline is:
\[
\begin{array}{c}
\text{universal weights}+\text{rational integral global series}\\
\Downarrow\\
\text{coprime parity products}=\text{reduced }A/B\\
\Downarrow\\
\text{algebraic roots}+\operatorname{Aut}(\Qbarl/\mathbf Q)
  \text{-stable weight strata}\\
\Downarrow\\
\mathbf Q\text{-factors}\to\mathbf Z\text{-factors}
\to\text{realization independence}.
\end{array}
\]
Keeping these as separate theorem boundaries prevents a later
formalization from smuggling arithmetic algebraicity into the definition
of a Weil sheaf.
\end{formalization}

